\setchapterpreamble[u]{\margintoc}
\chapter{The Four Fundamental Subspaces}
\labch{four-subspaces}

\sources{Strang, \emph{Introduction to Linear Algebra}, 6th ed.\ \cite{strang2023ila},
§3.1--3.5; MIT 18.06 \cite{ocw1806}, Lectures~6--11.}

Elimination answers a question about one right-hand side. The subspaces in this
chapter answer the question for all right-hand sides at once, and they do it
with four numbers: two dimensions in the input space and two in the output
space.

\section{Vector spaces and subspaces}
\labsec{subspaces}

\begin{definition}
A \emph{subspace} of \(\R^{n}\) is a non-empty set \(V\subseteq\R^{n}\) closed
under addition and scalar multiplication: if \(\vect{u},\vect{v}\in V\) and
\(c\in\R\) then \(\vect{u}+\vect{v}\in V\) and \(c\vect{v}\in V\).
\end{definition}

Taking \(c=0\) shows that every subspace contains \(\vect{0}\), which is the
quickest test to apply first. A plane through the origin in \(\R^{3}\) is a
subspace; the same plane shifted away from the origin is not, since it fails
that test.

\marginnote{Closure is the whole content of the definition. It says that the two
operations the subject is built from cannot carry you out of the set --- which
is exactly what makes linear questions answerable inside it.}

The subspaces that matter here are not chosen arbitrarily. Each of the four is
attached to a matrix and describes one aspect of what that matrix does.

\section{The column space}
\labsec{column-space}

The column space \(\Col(A)\subseteq\R^{m}\) was defined in
\refch{vectors-and-column-space} as the set of all combinations of the columns.
It is a subspace because a combination of combinations is again a combination,
and it settles solvability: \(A\vect{x}=\vect{b}\) has a solution precisely when
\(\vect{b}\in\Col(A)\).

Throughout this chapter I use
\[
	A=\begin{bmatrix}
		1&2&2&2\\
		2&4&6&8\\
		3&6&8&10
	\end{bmatrix}\in\R^{3\times4}.
\]
Its second column is twice its first, so the columns are dependent from the
start and the column space is smaller than \(\R^{3}\).

\section{The nullspace}
\labsec{nullspace}

\begin{definition}
The \emph{nullspace} of \(A\in\R^{m\times n}\) is
\[
	\Null(A)=\{\vect{x}\in\R^{n}: A\vect{x}=\vect{0}\}\subseteq\R^{n}.
\]
\end{definition}

It is a subspace: if \(A\vect{u}=\vect{0}\) and \(A\vect{v}=\vect{0}\) then
\(A(\vect{u}+c\vect{v})=\vect{0}\). Unlike the column space it lives in the
\emph{input} space, and it is exactly the object that \refsec{span-independence}
showed controls how many solutions a system has.

To compute it, eliminate. Subtracting \(2\times\)row~1 from row~2 and
\(3\times\)row~1 from row~3, then subtracting the new row~2 from the new row~3,
\[
	A\longrightarrow
	\begin{bmatrix}1&2&2&2\\0&0&2&4\\0&0&2&4\end{bmatrix}
	\longrightarrow
	U=\begin{bmatrix}1&2&2&2\\0&0&2&4\\0&0&0&0\end{bmatrix}.
\]

Two pivots appear, in columns~1 and~3. The row of zeros is not an accident: it
records that row~3 of \(A\) was the sum of rows~1 and~2 all along.

\begin{definition}
Columns containing a pivot are \emph{pivot columns}; the rest are \emph{free
columns}. The unknowns attached to free columns are \emph{free variables},
because they may be assigned any values whatever, after which the pivot
variables are determined.
\end{definition}

Here \(x_2\) and \(x_4\) are free. Setting \(x_2=1,x_4=0\) and back-substituting
gives \(x_3=0\) and \(x_1=-2\); setting \(x_2=0,x_4=1\) gives \(x_3=-2\) and
\(x_1=2\). So
\[
	\vect{s}_1=\begin{bmatrix}-2\\1\\0\\0\end{bmatrix},
	\qquad
	\vect{s}_2=\begin{bmatrix}2\\0\\-2\\1\end{bmatrix},
	\qquad
	\Null(A)=\operatorname{span}\{\vect{s}_1,\vect{s}_2\}.
\]

\marginnote{Check both by hand. \(A\vect{s}_1=-2\vect{a}_1+\vect{a}_2\) and the
second column is twice the first, so the result is \(\vect{0}\). Similarly
\(A\vect{s}_2=2\vect{a}_1-2\vect{a}_3+\vect{a}_4=\vect{0}\).}

\section{Rank, echelon form, and \texorpdfstring{\(A=CR\)}{A=CR}}
\labsec{rank-cr}

\begin{definition}
The \emph{rank} \(r\) of \(A\) is the number of pivots produced by elimination.
\end{definition}

For the running example \(r=2\). Continuing elimination until each pivot is~1
and the entries above it are cleared produces the \emph{reduced row echelon
form}
\[
	R_0=\begin{bmatrix}1&2&0&-2\\0&0&1&2\\0&0&0&0\end{bmatrix}.
\]

Discarding the zero row leaves \(R\in\R^{2\times4}\), and collecting the pivot
columns of the original matrix gives
\[
	C=\begin{bmatrix}1&2\\2&6\\3&8\end{bmatrix},
	\qquad
	R=\begin{bmatrix}1&2&0&-2\\0&0&1&2\end{bmatrix}.
\]

\begin{theorem}
Every matrix factors as \(A=CR\), where the \(r\) columns of \(C\) are the pivot
columns of \(A\) and the \(r\) rows of \(R\) are the non-zero rows of the reduced
row echelon form.
\end{theorem}

Multiplying out confirms it here: the first column of \(CR\) is
\(1\cdot\vect{c}_1=[1,2,3]\tp\), the second is \(2\vect{c}_1=[2,4,6]\tp\), the
third is \(\vect{c}_2=[2,6,8]\tp\), and the fourth is
\(-2\vect{c}_1+2\vect{c}_2=[2,8,10]\tp\).

\begin{remark}
\(A=CR\) is the cleanest statement of what rank means. The columns of \(A\) are
combinations of \(r\) independent columns, and the rows of \(A\) are
combinations of \(r\) independent rows --- the \emph{same} \(r\) in both cases.
Row rank and column rank being equal is often presented as a surprise; here it
is a consequence of a single factorisation.
\end{remark}

\section{The complete solution of \texorpdfstring{\(A\vect{x}=\vect{b}\)}{Ax=b}}
\labsec{complete-solution}

Take \(\vect{b}=[1,6,7]\tp\). Solving with the free variables set to zero gives a
\emph{particular} solution: the pivot variables must satisfy
\(x_1\vect{a}_1+x_3\vect{a}_3=\vect{b}\), that is \(x_1+2x_3=1\) and
\(2x_1+6x_3=6\), whence \(x_3=2\) and \(x_1=-3\); the third equation
\(3x_1+8x_3=-9+16=7\) is consistent. So
\[
	\vect{x}_p=\begin{bmatrix}-3\\0\\2\\0\end{bmatrix},
	\qquad
	\vect{x}=\vect{x}_p+c_1\vect{s}_1+c_2\vect{s}_2\ \ \text{for all }c_1,c_2\in\R.
\]

\begin{proposition}
If \(A\vect{x}=\vect{b}\) is solvable, its solution set is
\(\vect{x}_p+\Null(A)\): one particular solution shifted by the whole nullspace.
\end{proposition}

\begin{proof}
Any solution \(\vect{x}\) satisfies \(A(\vect{x}-\vect{x}_p)=\vect{b}-\vect{b}
=\vect{0}\), so \(\vect{x}-\vect{x}_p\in\Null(A)\). Conversely adding any
nullspace vector to \(\vect{x}_p\) leaves the product unchanged.
\end{proof}

\marginnote{This is the same statement as the trichotomy in
\refch{vectors-and-column-space}, now with the geometry supplied: the solution
set is a shifted copy of a subspace, so it is empty, a single point, or an
unbounded flat piece.}

\section{Independence, basis, and dimension}
\labsec{basis-dimension}

\begin{definition}
A \emph{basis} of a subspace \(V\) is a set of vectors that is linearly
independent and spans \(V\). The number of vectors in a basis is the
\emph{dimension} of \(V\), written \(\dim V\).
\end{definition}

That the number is the same for every basis is the fact that makes dimension
well defined. For the running example, \(\{\vect{s}_1,\vect{s}_2\}\) is a basis
of the nullspace and the two pivot columns form a basis of the column space, so
\[
	\dim\Col(A)=2,
	\qquad
	\dim\Null(A)=2 .
\]

\section{The four subspaces}
\labsec{four}

Applying the two constructions to \(A\) and to \(A\tp\) yields four subspaces in
total.

\begin{description}
	\item[Column space \(\Col(A)\subseteq\R^{m}\)] spanned by the columns;
	dimension \(r\).
	\item[Nullspace \(\Null(A)\subseteq\R^{n}\)] vectors killed by \(A\);
	dimension \(n-r\).
	\item[Row space \(\Col(A\tp)\subseteq\R^{n}\)] spanned by the rows;
	dimension \(r\).
	\item[Left nullspace \(\Null(A\tp)\subseteq\R^{m}\)] the vectors
	\(\vect{y}\) satisfying \(A\tp\vect{y}=\vect{0}\); dimension \(m-r\).
\end{description}

For the running example \(m=3\), \(n=4\), \(r=2\), so the dimensions are
\(2,2,2,1\). The left nullspace is spanned by \([1,1,-1]\tp\), which encodes the
dependency noticed during elimination: row~1 plus row~2 equals row~3.

\begin{theorem}[Rank--nullity]
For \(A\in\R^{m\times n}\) of rank \(r\),
\[
	\dim\Col(A)+\dim\Null(A)=r+(n-r)=n .
\]
\end{theorem}

\begin{proof}
Each of the \(n\) columns is either a pivot column or a free column. The pivot
columns form a basis of \(\Col(A)\), giving \(r\); each free column contributes
exactly one special solution, and those are independent because each carries a
\(1\) in a position where the others carry \(0\), giving \(n-r\).
\end{proof}

\begin{remark}
The counting has a plain reading. The input space \(\R^{n}\) splits its \(n\)
degrees of freedom into \(r\) that the matrix records and \(n-r\) that it
discards. Nothing else can happen, and no amount of cleverness in choosing
\(\vect{b}\) will change the split.
\end{remark}

\refch{orthogonality} sharpens this picture considerably: the two subspaces
inside \(\R^{n}\) are not merely complementary in dimension, they are orthogonal
complements, and likewise for the two inside \(\R^{m}\).
